\documentclass[11pt,letterpaper]{article}
\usepackage{amsmath,amssymb,amsthm}
\newtheorem{claim}{Claim}
\title{Reflections on Reading Slowly}
\author{Imogen Tarrant}
\date{December 1, 2017}
\begin{document}
\maketitle

\section{An Unfashionable Habit}
There is something faintly subversive about reading a single book for a
whole month. In an age of constant notification, sustained attention has
become an extraordinary accomplishment, and the pleasures of a
well-constructed paragraph are easily overlooked.

\begin{quotation}
\noindent To read slowly is not to read less; it is to allow each sentence to
finish what it has begun to say.
\end{quotation}

Several practical suggestions follow, drawn from the references
\cite{woolf, eco}.

\subsection{Guidelines}
\begin{enumerate}
\item Choose a book you suspect is difficult.
\item Read at a fixed hour, preferably at the same table.
\item Keep a notebook, and record
  \begin{itemize}
  \item unfamiliar words,
  \item passages that surprised you,
  \item questions for later.
  \end{itemize}
\item[$\star$] Never finish a chapter merely to finish it.
\end{enumerate}

\section{A Modest Mathematical Aside}
Suppose a reader finishes $p$ pages per day of a book of $N$ pages. The time
taken is $N/p$ days, so halving the speed doubles the time:
\[
\frac{N}{p/2} = 2\cdot\frac{N}{p}.
\]
\begin{claim}
A reader who reads $p$ pages daily for $d$ days has read $pd$ pages.
\end{claim}

\section*{Verbatim Diary Entry}
Here is a typical note, reproduced exactly:
\begin{verbatim}
Day 9: Chapter 4. Long sentence on p. 72.
       Read it three times. Still lovely.
\end{verbatim}
The macro \verb+\emph{slowly}+ would, of course, emphasise the word.

\begin{thebibliography}{9}
\bibitem{woolf} Virginia Woolf, \emph{How Should One Read a Book?}, 1932.
\bibitem{eco} Umberto Eco, \emph{On Literature}, Harcourt, 2004.
\end{thebibliography}
\end{document}
